\section{Uvod}
\label{sec:uvod}

Simetrična kriptografija predstavlja jednu od osnovnih tehnoloških
komponenti zaštite digitalnih komunikacija. Tokom poslednjih nekoliko
decenija menjali su se računarski kapaciteti, brzine komunikacionih
mreža, način implementacije kriptografskih algoritama i bezbednosni
zahtevi, dok je potreba za efikasnim šifrovanjem velikih količina
podataka ostala prisutna. Posebno značajan period u ovom razvoju bila
je tranzicija sa standarda Data Encryption Standard (DES), preko
njegovog višestrukog naslednika Triple DES odnosno Triple Data
Encryption Algorithm (TDEA), ka standardu Advanced Encryption Standard
(AES).

DES je tokom više decenija imao izuzetno značajnu ulogu u praktičnoj
kriptografiji, ali je njegova efektivna dužina ključa od 56 bita
vremenom postala njegovo najizraženije praktično ograničenje.
Problem nije proistekao iz potpunog kriptoanalitičkog rušenja njegove
unutrašnje konstrukcije, već pre svega iz činjenice da je prostor od
$2^{56}$ mogućih ključeva postao dostižan specijalizovanom hardveru.
Electronic Frontier Foundation (EFF)
demonstrirala je 1998. godine ovu činjenicu sistemom poznatim kao
\emph{Deep Crack}: DES ključ pronađen je iscrpnom pretragom za približno
56 časova, uz prosečnu brzinu od oko $8.8\times10^{10}$ ispitanih
ključeva u sekundi~\cite{eff-des-cracker-1998}. Ovakvi praktični
rezultati pokazali su da bezbednost sistema ne može biti procenjivana
samo na osnovu složenosti same šifre, već mora uzeti u obzir i veličinu
prostora ključeva, cenu računarskih resursa i razvoj specijalizovanog
hardvera.

TDEA je produžio upotrebljivost DES infrastrukture, uz veći računarski
trošak i nasleđeno ograničenje bloka. Oba standarda kasnije su povučena
za novu zaštitu u NIST okviru~\cite{nist-des-withdrawal-2005,
nist-tdea-withdrawal-2023}. Otvoreni AES proces doveo je do izbora
Rijndaela i standardizacije u FIPS~197, sa ciljem dugoročne bezbednosti
i efikasne primene na različitim platformama
~\cite{nist-aes-report-2001,nist-aes-development-2021,fips197-2001}.

Period oko 2001--2003. godine zato predstavlja posebno interesantnu
tačku za retrospektivnu analizu. U tom trenutku DES je još postojao u
brojnim realnim sistemima, Triple DES je predstavljao praktičnu
migracionu mogućnost, dok je AES bio novi standard čija dugoročna
bezbednost, zastupljenost i implementaciona svojstva tek trebalo da
budu potvrđena višegodišnjom praktičnom primenom. Dve decenije kasnije
moguće je analizirati koliko su tadašnja očekivanja bila opravdana i
koji problemi, tada manje vidljivi, danas predstavljaju važan deo
diskusije o simetričnoj kriptografiji.

Ovaj rad je motivisan upravo tim tehnološkim periodom tranzicije.
Njegov cilj nije reprodukcija ranijeg magistarskog rada autora iz
2003. godine, već nova retrospektivna analiza razvoja od DES-a preko
TDEA do AES-a, zasnovana na danas dostupnoj istorijskoj
dokumentaciji, standardima, naučnoj literaturi i eksperimentalnim
merenjenjima. Centralno pitanje rada jeste kako su se očekivanja
formirana u vreme izbora i početnog uvođenja AES-a pokazala nakon više
od dve decenije njegove praktične upotrebe.


\subsection{Istorijski trenutak i motivacija}
\label{subsec:istorijski_trenutak}

Početkom novog milenijuma novi standard morao je da objedini veću
sigurnosnu rezervu i efikasnost potrebnu komunikacionim sistemima.
Dokumentacija AES procesa omogućava da se ti kriterijumi rekonstruišu
nezavisno od kasnijeg uspeha Rijndaela~\cite{nist-aes-report-2001}.
Period od približno četvrt veka zatim pruža osnovu za poređenje tih
očekivanja sa razvojem standarda, protokola i implementacija.

\subsection{Istraživačka pitanja i doprinos}
\label{subsec:istrazivacka_pitanja}

Rad je koncipiran kao pregledno-eksperimentalna retrospektivna studija.
Ne predlaže se novi simetrični algoritam niti nova kriptoanalitička
metoda. Polazište je istorijska tranzicija od DES-a preko TDEA ka AES-u,
ali je centralni cilj širi: proveriti kako su se dokumentovana očekivanja
iz perioda izbora AES-a pokazala nakon približno četvrt veka razvoja
kriptografije, procesorskih arhitektura, komunikacionih protokola i
praktičnih implementacija.

Analiza je organizovana oko sledećih istraživačkih pitanja:

\begin{enumerate}
    \item[\textbf{IP1}] Koja su bezbednosna, implementaciona i
    arhitektonska ograničenja DES-a i TDEA dovela do potrebe za novim
    standardom?

    \item[\textbf{IP2}] Koji su zahtevi i očekivanja postavljeni pred
    AES tokom procesa izbora 1997--2001. godine i zbog kojih osobina je
    Rijndael izabran?

    \item[\textbf{IP3}] U kojoj meri su se očekivanja u pogledu
    bezbednosti, performansi, implementacione fleksibilnosti i
    dugovečnosti AES-a potvrdila do 2026. godine?

    \item[\textbf{IP4}] Kako savremena procesorska i GPU akceleracija
    menjaju praktičnu cenu AES-a i koliko zaključak zavisi od toga da li
    se meri izolovana transformacija, device-resident izvršavanje ili
    kompletna end-to-end putanja?

    \item[\textbf{IP5}] Koja ograničenja i otvorena pitanja postaju
    vidljiva iz današnje perspektive, posebno u vezi sa autentifikovanom
    enkripcijom, bočnim kanalima, veličinom bloka, lightweight
    okruženjima i budućim kvantnim modelom napada?
\end{enumerate}

Prvi doprinos rada je pregledno-sintetički. Dokumentovani razlozi za
slabljenje DES-a, produženu tranzicionu ulogu TDEA i izbor Rijndaela
povezuju se sa kasnijim razvojem standarda i protokola do 2026. godine.
Pri tome se izričito razdvajaju očekivanja dokumentovana tokom AES procesa
od kasnijih tehnoloških razvoja, kako se AES-NI, GPU izvršavanje ili
savremena AEAD praksa ne bi retroaktivno pripisivali originalnom procesu
izbora.

Drugi i centralni eksperimentalni doprinos je metodološki. Uspostavljen
je validiran i ponovljiv okvir za mikrobenchmarking kriptografskih
implementacija sa funkcionalnom validacijom pre merenja, eksplicitnim
granicama tajmera, dokumentovanim seedovima i sesijama, čuvanjem raw,
summary, validation i metadata artefakata i agregacijom rezultata između
nezavisnih procesa.

Posebno mesto ima upareni protokol za poređenje CPU AES-NI i CUDA
AES-128 implementacija. Isti ulazi koriste se kroz poređene putanje, dok
se odvojeno posmatraju CPU transform-only, GPU resident i GPU end-to-end
obuhvat. Time se neposredno ispituje koliko zaključak o performansama
zavisi od merne granice. Takvo razdvajanje je važno jer visoka propusnost
izolovanog GPU kernela ne mora predstavljati propusnost praktične
host--device aplikacione putanje.

Treći doprinos predstavlja empirijska evaluacija ovako definisanog
okvira na dokumentovanoj savremenoj platformi. Ona obuhvata istorijsko
CPU poređenje DES/TDEA/AES, softverski AES naspram AES-NI, AES-GCM,
CUDA AES-128 baseline i upareno CPU--GPU poređenje, kao i ilustrativnu
pretragu redukovanog prostora ključeva i avalanche analizu. Ti rezultati
ne predstavljaju univerzalne performanse algoritama niti rekonstrukciju
hardvera iz vremena standardizacije AES-a. Njihova uloga je da
kvantifikuju razliku između primitive, implementacione putanje i
kompletnog sistema na jednoj reproducibilno opisanoj platformi.

Važan deo metodološkog doprinosa je i eksplicitno dokumentovanje
ograničenja i negativnih rezultata. CUDA kandidat lake optimizacije nije
doneo konzistentno poboljšanje, a samostalni baseline i uparena studija
nisu pokazali isti obrazac performansi na najvećim baferima. Takvi nalazi
se ne uklanjaju niti generalizuju bez dodatne kontrole stanja uređaja.
Ovim se performansni rezultat vezuje za protokol kojim je dobijen, umesto
za pretpostavljenu univerzalnu osobinu CPU-a ili GPU-a.

Eksperimentalni deo zato nije zamišljen kao potraga za novim rekordom
propusnosti, već kao kontrolisana demonstracija koliko se zaključak o
kriptografskim performansama menja kada se promeni implementaciona
putanja ili obuhvat merenja. Time eksperiment neposredno podržava
retrospektivnu temu rada: algoritam standardizovan početkom 2000-ih
ostaje isti, dok se način njegovog efikasnog i bezbednog izvršavanja
značajno menja.


\subsection{Metod izbora literature i vremenski obuhvat}
\label{subsec:metod_literature}

Pregled literature u ovom radu nije zamišljen kao formalni sistematski
pregled u smislu medicinskih ili meta-analitičkih protokola, već kao
strukturisan istorijsko-tehnički pregled. Izvori se biraju hijerarhijski,
u zavisnosti od vrste tvrdnje koja se proverava.

Za definicije algoritama, parametre standarda, datume standardizacije i
status kriptografskih preporuka prioritet imaju primarni normativni
izvori, pre svega publikacije NIST-a, Federal Information Processing
Standards (FIPS), NIST Special Publications i relevantni dokumenti
Internet Engineering Task Force (IETF). Istorijski deo razvoja AES-a
zasniva se prvenstveno na dokumentaciji samog AES procesa i izveštajima
NIST-a iz perioda 1997--2001. godine
~\cite{nist-aes-report-2001,fips197-2001}.

Za pitanja kriptoanalize, napada putem bočnih kanala, implementacionih
slabosti i performansi koriste se originalni ili recenzirani naučni
radovi kad god su dostupni. Sekundarni izvori, stručne diskusije i
tehnički članci mogu se koristiti za identifikovanje relevantnih
pitanja, ali se ključni zaključci ne zasnivaju na njima kada postoji
primarni standard ili originalna naučna publikacija.

Posebna pažnja posvećena je vremenskom kontekstu izvora. Tvrdnje o
očekivanjima od AES-a zasnivaju se na dokumentima nastalim pre ili
tokom njegovog izbora i standardizacije, kako kasnija saznanja ne bi
bila retroaktivno pripisana učesnicima tog procesa. Nasuprot tome,
procena ishoda koristi kasnije izvore do vremenskog preseka 2026.
godine. Na taj način se razdvajaju dve perspektive:
\emph{šta se tada očekivalo} i \emph{šta je kasnije zabeleženo u
praksi i standardima}.

Kada ista specifikacija postoji u više revizija, istorijska analiza
koristi izdanje relevantno za posmatrani period, dok se za aktuelni
status koristi najnovija važeća ili poslednja povučena verzija.
Tako se istorijsko i ažurirano izdanje FIPS~197 koriste prema ulozi
tvrdnje u analizi~\cite{fips197-2023}.

Vremenski obuhvat rada zato počinje razvojem i slabljenjem praktične
pozicije DES-a, sa posebnim fokusom na period od druge polovine
devedesetih godina, prati izbor i standardizaciju AES-a do 2001.
godine i zatim razvoj njegove primene i implementacije do 2026.
godine.


\subsection{Telekomunikacioni kontekst i organizacija rada}
\label{subsec:telekom_kontekst}

Telekomunikacioni sistemi predstavljaju posebno pogodno okruženje za
analizu evolucije simetričnih šifara, jer se u njima bezbednosni zahtevi
ne mogu posmatrati odvojeno od performansi. Šifrovanje podataka u
prenosu mora biti dovoljno brzo da ne ograničava propusnost veze, dok
latencija, procesorsko opterećenje, potrošnja energije i dostupni
hardverski resursi mogu biti jednako značajni kao teorijska sigurnost
algoritma.

Istovremeno, komunikaciona infrastruktura često ima dug životni vek.
Kriptografski mehanizam koji je u jednom trenutku savremen može ostati
prisutan godinama zbog kompatibilnosti sa postojećom opremom,
protokolima i implementacijama. Upravo je takav problem relevantan za
prelaz sa DES-a preko TDEA na AES: bezbednosni standard može
biti zamenjen relativno brzo na normativnom nivou, dok praktična
migracija infrastrukture traje znatno duže.

Razvoj mrežnih protokola takođe pokazuje da značaj AES-a nije ograničen
na samu blokovsku šifru. Način na koji se šifra koristi postao je
jednako važan kao izbor primitive. RFC~4106 je još 2005. godine
definisao AES u Galois/Counter Mode (AES-GCM) za IPsec ESP radi
istovremenog obezbeđivanja poverljivosti i autentikacije porekla
podataka. U dokumentu je posebno naglašena mogućnost efikasne
hardverske implementacije pri brzinama od $10$~Gbit/s i više
~\cite{rfc4106}.

Kontinuitet ove uloge može se videti i u savremenim protokolima.
Aktuelna specifikacija TLS~1.3, RFC~9846 iz 2026. godine, zahteva
implementaciju skupa \texttt{TLS\_\allowbreak AES\_\allowbreak 128\_\allowbreak GCM\_\allowbreak SHA256}, dok
\texttt{TLS\_\allowbreak AES\_\allowbreak 256\_\allowbreak GCM\_\allowbreak SHA384} navodi kao preporučenu
implementaciju~\cite{rfc9846}. To ne znači da je AES jedina savremena
simetrična šifra niti da je jednako zastupljen u svim
telekomunikacionim sistemima, ali pokazuje njegovu kontinuiranu
relevantnost više od dve decenije nakon standardizacije.

Razlika između primitive, načina rada i protokola uspostavlja se u
Poglavlju~\ref{sec:osnove} i primenjuje kroz ostatak analize.

Nakon uvodnog poglavlja, u Poglavlju~\ref{sec:osnove} daju se osnovni
pojmovi simetrične kriptografije potrebni za dalje poređenje.
Zatim se analiziraju DES i TDEA, njihove konstrukcije i razlozi
za postepeno napuštanje. Posebno poglavlje posvećeno je procesu izbora
AES-a, nakon čega sledi prikaz njegove konstrukcije i implementacionih
bezbednosnih pitanja. Na toj osnovi vrši se neposredno poređenje tri
generacije standarda.

Drugi deo rada razmatra primene AES-a u telekomunikacionim i mrežnim
sistemima. Eksperimentalni deo zatim ispituje performanse konkretnih
softverskih implementacija, efekat savremene hardverske akceleracije,
ilustrativnu iscrpnu pretragu redukovanog prostora ključeva i
avalanche efekat. Rezultati pregleda i eksperimenata objedinjuju se u
retrospektivnom poglavlju koje poredi očekivanja iz perioda
2001--2003. sa stanjem do 2026. godine. Završni deo rada razmatra
ograničenja, savremene izazove i pravce mogućeg budućeg istraživanja.
